\documentclass{article}

\usepackage{amsmath,amssymb}
\usepackage{xurl}
\usepackage[hidelinks]{hyperref}
\setlength{\emergencystretch}{2em}

% Shared commands for notes in docs/paper-gaps/.
% Notation follows the LDT paper (references/ldt-paper/).

\newcommand{\Lean}{\textsc{Lean}}
\newcommand{\leanid}[1]{\textsf{#1}}
\newcommand{\ghissue}[1]{\href{https://github.com/LionSR/MIPStarRE/issues/#1}{GitHub issue \##1}}
\newcommand{\ghpr}[1]{\href{https://github.com/LionSR/MIPStarRE/pull/#1}{GitHub PR \##1}}

% --- Paper-compatible notation ---
\newcommand{\F}{\mathbb{F}}
\newcommand{\N}{\mathbb{N}}
\newcommand{\C}{\mathbb{C}}
\newcommand{\tr}{\operatorname{tr}}
\newcommand{\ot}{\otimes}
\newcommand{\E}{\mathop{\bf E\/}}
\newcommand{\bx}{\mathbf{x}}
\newcommand{\by}{\mathbf{y}}
\newcommand{\bu}{\mathbf{u}}
\newcommand{\bv}{\mathbf{v}}

% Approximation relation (paper uses \approx_\eps and \simeq_\zeta)
\newcommand{\approxeps}{\approx_{\varepsilon}}
\newcommand{\simgeq}{\simeq_{\zeta}}

% Polynomial spaces (paper uses \polyfunc, \polysub, \polymeas)
\newcommand{\polyfunc}[3]{\operatorname{Poly}(#1,#2,#3)}
\newcommand{\polysub}[3]{\operatorname{PolySub}(#1,#2,#3)}
\newcommand{\polymeas}[3]{\operatorname{PolyMeas}(#1,#2,#3)}


\title{The Restricted Interface of the Main Formal Soundness Theorem}
\author{}
\date{2026-05-01}

\begin{document}
\maketitle

\section{The assertion}

The main formal soundness theorem in \cite[Section~\emph{The test}]{Ji2020LowIndividualDegree}
starts from a general projective strategy for the $(m,q,d)$ low individual degree
test. In the paper, such a strategy consists of a bipartite state on
$\mathcal H_{\mathrm A}\otimes\mathcal H_{\mathrm B}$ and separate projective
measurement families for the two provers. The two local Hilbert spaces are not
assumed to be the same, and the starting state is not assumed to be symmetric.
The test definition takes $m$ and $d$ to be non-negative integers, and the
main theorem as printed does not impose a separate hypothesis $d>0$.\footnote{This
note is part of the retrospective audit in \ghissue{930}. The paper passages are
\path{references/ldt-paper/test_definition.tex}, lines~10--14, 98--115, and
180--202, together with the reduction in
\path{references/ldt-paper/inductive_step.tex}, lines~26--66. The previously
documented large-$k$ correction is separate; see
\path{docs/paper-gaps/issue-906-main-formal-k-bound.tex}.}

After stating the theorem for this general strategy, the paper symmetrizes it by
adding a two-dimensional role register on each side. In the proof, the paper
first treats the case in which the two local Hilbert spaces have the same
dimension and says that the construction extends to different dimensions. Even
in that simplified proof setting, however, the original state is not assumed to
be swap-invariant. The symmetrized state is a sum of two role sectors: one sector
carries the original state, and the other sector carries the state with the
tensor factors interchanged. The purpose of this step is precisely to pass from
a not-necessarily-symmetric general strategy to a symmetric strategy to which the
main induction theorem can be applied.

\section{The formal interface}

The present formal development now contains a paper-faithful two-space container
for general projective strategies, but the public main-formal interface has not
yet been generalized to it. The current theorem is stated for the same-space
special case in which both provers use a common local carrier and the starting
bipartite state is equipped with swap-invariance data. The role-register
symmetrization bridge used by the theorem has the same restricted input shape:
it builds the symmetrized strategy only from this same-space, already
swap-compatible strategy object.\footnote{The two-space strategy container is
\leanid{MIPStarRE.LDT.ProjStrat}; the same-space special case is
\leanid{MIPStarRE.LDT.SameSpaceProjStrat}. Pull request \ghpr{958} closed
\ghissue{560} by making \leanid{ProjStrat} two-space, but
\leanid{MIPStarRE.LDT.Test.mainFormal} in
\doclink{MIPStarRE/LDT/Test/MainTheorem.lean} and the Step~1 bridge
\leanid{MIPStarRE.LDT.SameSpaceProjStrat.strategySymmetrization} in
\doclink{MIPStarRE/LDT/Test/SymmetrizationBridge.lean} still take
\leanid{SameSpaceProjStrat}.}

The blueprint records the same interface boundary. It describes the paper's
general projective strategy while noting that the formal strategy layer separates
the two-space container from the same-space surrogate used by the present
role-register symmetrization pipeline. Its current main-formal statement also
includes the additional hypothesis $d>0$.\footnote{See
\path{blueprint/src/chapter/ch02_test.tex}, lines~35--43 for the strategy split
and lines~102--122 for the main-formal blueprint statement.}

Thus the formal theorem is not presently a theorem about every projective
strategy in the sense of the paper. A general strategy on
$\mathcal H_{\mathrm A}\otimes\mathcal H_{\mathrm B}$ with different local
spaces, or with a nonsymmetric starting state, is an input to the paper theorem
but not an input to the current public formal statement. To reach the paper
statement, the role-register construction must be carried out directly for the
two-space strategy container, producing a symmetric strategy on the appropriate
direct-sum role space without first assuming a common local carrier or
swap-invariance of the original state.

The formal theorem also assumes $d>0$. This is another statement-level
restriction absent from the printed theorem. It is harmless for the nontrivial
regime used by the current proof route, but it does not by itself establish the
constant-degree case $d=0$ covered by the paper's non-negative-degree wording.
A paper-exact public wrapper would need either to remove this hypothesis from
the formal pipeline or to add a separate argument for $d=0$.

\section{Conclusion}

The discrepancy is a formal-interface restriction, not a flaw in the paper's
strategy-reduction argument. The paper theorem starts from an arbitrary general
projective strategy and then symmetrizes it. The current formal statement starts
from a narrower same-space, swap-compatible strategy object, and it also assumes
positive degree. These restrictions should remain visible until the Step~1
symmetrization bridge and the main-formal theorem are generalized to the
paper-faithful two-space strategy interface, or until separate wrappers prove
that the omitted cases reduce to the formal statement.

\bibliographystyle{alpha}
\bibliography{references}

\end{document}
